\documentclass[main.tex]{subfiles}


\begin{document}

%from pagenumber to up
% \phantomsection 
% \hypertarget{toc}{}

 

 

% номер секции вручную
\setcounter{section}{2
}
\addtocounter{section}{-1} 

\section{Заряд}


Заряженные тела могут как притягиваться друг к другу, так и отталкиваться. В общем случае характер взаимодействия таких тел определяется их зарядами.

\textbf{Заряд} ($q$ [Кл])~--- это характеристика способности заряженного тела взаимодействовать с другим заряженным телом. Чем больше заряд данного тела, тем с большей силой это тело притягивается (или отталкивается) к другому заряженному телу при прочих равных условиях. Положительно заряженное тело имеет положительный заряд, отрицательно заряженное~--- отрицательный заряд. Для краткости всякое заряженное тело или частицу называют просто зарядом.

На рис.\,\ref{fig:p89} показаны два заряженных тела.

    \begin{figure}[H]
    \centering

    \begin{tikzpicture}[font=\small,            line cap=round,                             >={Stealth[inset=0pt,round,length=3mm, angle'=20]},vec/.style={>={Stealth[inset=0pt,round,length=3.5mm, angle'=20]}}]


\tikzset{atom/.pic={
  \begin{scope}[shift={({rnd*360}:.05)}]
  \shade[ball color=red] (0,0) circle (\dn);
  \shade[ball color=NavyBlue!70,rotate={rnd*360}] ([shift={(180:{\dn cm})}]0,0) circle (\de);
  \end{scope}}
}

\tikzset{atompo/.pic={
  \begin{scope}[shift={({rnd*360}:.05)}]
  \shade[ball color=red] (0,0) circle (\dn);
  \end{scope}}
}

\tikzset{atomne/.pic={
  \begin{scope}[shift={({rnd*360}:.05)},rotate={rnd*360}]
  \shade[ball color=red] (0,0) circle (\dn);
  \shade[ball color=NavyBlue!70] ([shift={(180:{\dn cm})}]0,0) circle (\de);
  \shade[ball color=NavyBlue!70] ([shift={(0:{\dn cm})}]0,0) circle (\de);
  \end{scope}}
}

\def\dn{.15}
\def\de{.1}


%solid
\fill[red,nearly transparent] (0,0) rectangle++ (4,2);
\draw[] (0,0) rectangle++ (4,2);

\fill[NavyBlue,nearly transparent] (7,0) rectangle++ (4,2);
\draw[] (7,0) rectangle++ (4,2);


%T
\node[above] at (2,2) {$q_\text{А}$};
\node[above] at (9,2) {$q_\text{Б}$};


%q
\node[below] at (2,0) {А};
\node[below] at (9,0) {Б};



%atoms left
\foreach \y in {0,...,2}
\foreach \x in {0,...,4}
\pic[shift={(.4,.4)}] at ({0+(4-.8)/5*\x},{0+(2-.8)/2*\y}) {atom};

\foreach \y in {0,...,2}
\pic[shift={(.4,.4)}] at ({0+(4-.8)/5*5},{0+(2-.8)/2*\y}) {atompo};

%%%

%atoms right
\foreach \y in {0,...,2}
\foreach \x in {1,...,5}
\pic[shift={(7.4,.4)}] at ({0+(4-.8)/5*\x},{0+(2-.8)/2*\y}) {atom};

\foreach \y in {0,...,2}
\pic[shift={(7.4,.4)}] at ({0+(4-.8)/5*0},{0+(2-.8)/2*\y}) {atomne};



    \end{tikzpicture}
\caption{Заряженные тела}
\label{fig:p89}
\end{figure}

Каждое тело на рис.\,\ref{fig:p89} упрощенно составлено из протонов (красные шары) и электронов (синие шары)\footnote{Внимание уделено именно заряженным частицам, находящимся в атомах тела. Разница в размерах шаров на рисунке~\ref{fig:p89} подчеркивает различие масс соответствующих частиц: масса протона почти в 2000~раз больше массы электрона.}. Тело~А имеет положительный заряд~$q_\text{А}>0$: в нем протоны преобладают над электронами. Наоборот, тело~Б несет отрицательный заряд~$q_\text{Б}<0$: в нем уже электроны преобладают над протонами. 

% Приблизительно такая картина распределения зарядов в телах получается после трения, например, пластикового стержня о шерсть.  

\emph{Заряды (заряженные тела) разных знаков притягиваются друг к другу, а заряды одного знака друг от друга отталкиваются.} Сказанное иллюстрируется рисунком~\ref{fig:p90} (красным шарам на нитях сообщены положительные заряды, синим~--- отрицательные).

    \begin{figure}[H]
    \centering

    \begin{tikzpicture}[font=\small,            line cap=round,                             >={Stealth[inset=0pt,round,length=3mm, angle'=20]},vec/.style={>={Stealth[inset=0pt,round,length=3.5mm, angle'=20]}}]


\filldraw[lightgray,semitransparent] (0,3) rectangle++ (3.1,.3);  
\draw (0,3) --++ (3.1,0);


%1
\draw (.6,3) --++ (-75:2);
\path (.6,3) --++ (-75:2.15) coordinate (cir);
\shade[ball color=red] (cir) coordinate (C) circle (.15);

\draw (2.5,3) --++ (-105:2);
\path (2.5,3) --++ (-105:2.15) coordinate (cir);
\shade[ball color=NavyBlue] (cir) coordinate (C) circle (.15);


%2
\begin{scope}[xshift=4.5cm]

\filldraw[lightgray,semitransparent] (0,3) rectangle++ (3.1,.3);  
\draw (0,3) --++ (3.1,0);


\draw ({.6+cos(75)*2.15},3) --++ (-105:2);
\path ({.6+cos(75)*2.15},3) --++ (-105:2.15) coordinate (cir);
\shade[ball color=red] (cir) coordinate (C) circle (.15);

\draw ({2.5-cos(75)*2.15},3) --++ (-75:2);
\path ({2.5-cos(75)*2.15},3) --++ (-75:2.15) coordinate (cir);
\shade[ball color=red] (cir) coordinate (C) circle (.15);
    
\end{scope}


%3
\begin{scope}[xshift=9cm]

\filldraw[lightgray,semitransparent] (0,3) rectangle++ (3.1,.3);  
\draw (0,3) --++ (3.1,0);


\draw ({.6+cos(75)*2.15},3) --++ (-105:2);
\path ({.6+cos(75)*2.15},3) --++ (-105:2.15) coordinate (cir);
\shade[ball color=NavyBlue] (cir) coordinate (C) circle (.15);

\draw ({2.5-cos(75)*2.15},3) --++ (-75:2);
\path ({2.5-cos(75)*2.15},3) --++ (-75:2.15) coordinate (cir);
\shade[ball color=NavyBlue] (cir) coordinate (C) circle (.15);
    
\end{scope}


    \end{tikzpicture}
\caption{Притяжение и отталкивание зарядов}
\label{fig:p90}
\end{figure}

\textbf{Элементарный электрический заряд} ($e$ [Кл])~--- это величина наименьшего заряда, существующего в природе ($e\approx1{,}6\cdot10^{-19}$~Кл). Протон и электрон являются частицами с минимальной возможной величиной заряда. Заряд протона равен $q_p\approx1{,}6\cdot10^{-19}$~Кл; заряд электрона равен $q_e\approx-1{,}6\cdot10^{-19}$~Кл, то есть его заряд равен заряду протона со знаком минус. 

Заряд тела обусловлен преобладанием в нем протонов или электронов. Если заряд тела $q>0$, то $q=Nq_p$ (где $N$~--- количество избыточных протонов). Если заряд тела $q<0$, то $q=Nq_e$ (где $N$~--- количество избыточных электронов).























 
\end{document}